% Options for packages loaded elsewhere
\PassOptionsToPackage{unicode}{hyperref}
\PassOptionsToPackage{hyphens}{url}
\PassOptionsToPackage{dvipsnames,svgnames,x11names}{xcolor}
%
\documentclass[
  11pt,
]{article}
\usepackage{amsmath,amssymb}
\usepackage{iftex}
\ifPDFTeX
  \usepackage[T1]{fontenc}
  \usepackage[utf8]{inputenc}
  \usepackage{textcomp} % provide euro and other symbols
\else % if luatex or xetex
  \usepackage{unicode-math} % this also loads fontspec
  \defaultfontfeatures{Scale=MatchLowercase}
  \defaultfontfeatures[\rmfamily]{Ligatures=TeX,Scale=1}
\fi
\usepackage{lmodern}
\ifPDFTeX\else
  % xetex/luatex font selection
\fi
% Use upquote if available, for straight quotes in verbatim environments
\IfFileExists{upquote.sty}{\usepackage{upquote}}{}
\IfFileExists{microtype.sty}{% use microtype if available
  \usepackage[]{microtype}
  \UseMicrotypeSet[protrusion]{basicmath} % disable protrusion for tt fonts
}{}
\makeatletter
\@ifundefined{KOMAClassName}{% if non-KOMA class
  \IfFileExists{parskip.sty}{%
    \usepackage{parskip}
  }{% else
    \setlength{\parindent}{0pt}
    \setlength{\parskip}{6pt plus 2pt minus 1pt}}
}{% if KOMA class
  \KOMAoptions{parskip=half}}
\makeatother
\usepackage{xcolor}
\usepackage[margin=1in]{geometry}
\setlength{\emergencystretch}{3em} % prevent overfull lines
\providecommand{\tightlist}{%
  \setlength{\itemsep}{0pt}\setlength{\parskip}{0pt}}
\setcounter{secnumdepth}{-\maxdimen} % remove section numbering
\ifLuaTeX
  \usepackage{selnolig}  % disable illegal ligatures
\fi
\usepackage{bookmark}
\IfFileExists{xurl.sty}{\usepackage{xurl}}{} % add URL line breaks if available
\urlstyle{same}
\hypersetup{
  colorlinks=true,
  linkcolor={Maroon},
  filecolor={Maroon},
  citecolor={Blue},
  urlcolor={Blue},
  pdfcreator={LaTeX via pandoc}}

\author{}
\date{}

\begin{document}

\section{P15.2: common full-mass domain and minimum full-form operator
data}\label{p15.2-common-full-mass-domain-and-minimum-full-form-operator-data}

Xean research manuscript, 4 October 2026. The complete proof passed an
independent review at its stated scope. Original P15.2 remains open
pending a separate audit of the full question. Energy integrability
specifies a common operator subdomain.

The new work is the scalar-deficit common-domain theorem, bare-current
reconstruction, extension of the reference kernel to all full-mass
inputs, quotient continuity and the explicit infinite-dimensional
geometric-kernel example. Published foundations are cited with
hypotheses and checked applications. Earlier supplied author/review
verdicts are evidence to reconstruct, not theorem grants.

\section{n1. Common full-mass and joint-linear
domains}\label{n1.-common-full-mass-and-joint-linear-domains}

\textbf{Scope.} This is a partial P15.2 result, independently
reconstructed from the stated literature foundations. No earlier proof
or review verdict is an assumption. Only P15.2 is investigated. The
energy integrability condition below specifies a useful operator domain;
no P15.1 energy theorem is claimed.

Let \(X\) be a connected smooth compact Kahler manifold of complex
dimension \(n\geq1\), let \(\alpha\) be pseudoeffective with
\(\operatorname{vol}(\alpha)=0\), let \(\theta\) be a smooth closed real
representative, and let \(\psi\in\operatorname{PSH}(X,\theta)\) not be
identically minus infinity. Set \[
d^c=\frac{i}{4\pi}(\bar\partial-\partial),\qquad dd^c=\frac{i}{2\pi}\partial\bar\partial,
\qquad T_u=\theta+dd^cu,\qquad P=T_\psi.
\] The retained type is the actual two-sided bounded comparison type \[
\mathcal{C}_\psi=\{u\in\operatorname{PSH}(X,\theta):\psi-C\leq u\leq\psi+C\text{ for some }C<\infty\}.
\] Comparisons are extended-valued global comparisons, not merely
comparisons on intersecting finite loci. A more-singular potential means
\(u\leq\psi+C\). Differences in integrals are evaluated on the
simultaneous finite locus and assigned zero elsewhere. Every measure
below gives zero mass to these pluripolar exceptional sets.

The essential standard foundations used here are the following, with
complete supplied printed bodies in the source annex.

\begin{itemize}
\tightlist
\item
  Boucksom--Eyssidieux--Guedj--Zeriahi (BEGZ), \emph{Monge-Ampere
  equations in big cohomology classes}, Acta Math. 205 (2010), 199--262,
  arXiv:0812.3674v2, Definition 1.1, Propositions 1.4 and 1.6, Theorem
  1.8: on compact Kahler \(X\), nonpluripolar products of at most \(n\)
  closed positive \((1,1)\) currents are well-defined finite closed
  positive currents, symmetric, multilinear on positive sums and
  plurifinely local; their trace measures do not charge pluripolar sets.
  Smooth factors commute with the construction. The empty product is
  \(1\). These are actual-current properties, not multilinearity of
  class positive products. All currents used below meet these
  hypotheses.
\item
  Darvas--Di Nezza--Lu (DDNL), \emph{Monotonicity of nonpluripolar
  products and complex Monge-Ampere equations with prescribed
  singularity}, Anal. PDE 11 (2018), 2049--2087, Proposition 2.1: for
  \(n\) smooth closed real backgrounds \(\theta_j\) with pseudoeffective
  classes and \(u_j,v_j\in\operatorname{PSH}(X,\theta_j)\) with the same
  bounded singularity type in each slot, the two top mixed masses are
  equal. Theorem 2.4: if \(u_j\) is less singular than \(v_j\) in each
  slot, its top mixed mass is at least the other mass. No big-class or
  small-unbounded-locus condition is present. Every application below
  uses precisely \(k\) potential slots in \(\theta\) and \(n-k\) smooth
  Kahler slots with zero potentials; the latter classes are
  pseudoeffective, and both comparisons in those slots are identical.
\item
  BEGZ Sections 1.3--1.4 and Definition 1.17: minimal-singularity
  currents exist in a pseudoeffective class, are less singular than
  every current in that class, and their top mass in a big class is its
  volume. At the boundary the volume is the decreasing limit along
  addition of a Kahler class. The positive product of nef classes agrees
  with cup product. We use the actual-current upper bound supplied by
  DDNL comparison with the minimal current, not an equality of arbitrary
  mass with class volume.
\item
  Duc-Viet Vu, \emph{Convexity of the class of currents with finite
  relative energy}, arXiv:2005.13241v3 (25 March 2022), Section 3: with
  pseudoeffective classes, positive reference currents \(P_j\), and a
  closed positive relative current \(R\), the relative full-mass tuple
  space consists of currents more singular than the references and with
  equal cohomology classes of their relative nonpluripolar products.
  Joint \(\chi\) energy is the sum over subsets \(J\) of the integrals
  of \(-\chi(\sum_j(u_j-\varphi_j))\) against the mixed relative
  products choosing \(T_j\) on \(J\) and \(P_j\) elsewhere. Here
  \(\chi\) is convex increasing and \(\chi(-\infty)=-\infty\). Theorem
  3.4 states convexity of the diagonal full-mass and diagonal
  joint-weight spaces. Corollary 3.2 states equality of the joint-weight
  spaces after a bounded-equivalent cohomologous reference change. The
  printed proofs invoke other published arguments; these two results are
  used as precisely scoped literature foundations, not as newly proved
  lemmas. We apply them only for \(d\geq1\), \(d\) equal slots \(P\),
  \(R=\omega^{n-d}\) smooth closed positive, and \(\chi(t)=t\). The
  products have top degree, so on connected compact oriented \(X\) their
  cohomology equality is exactly equality of masses: integration
  identifies \(H^{2n}(X,\mathbb{R})\) with \(\mathbb{R}\).
\end{itemize}

DDNL uses \(i\partial\bar\partial\), so its potential for the same
current is \(u/(2\pi)\). Vu uses
\(dd^c_s=i\partial\bar\partial/\pi=2dd^c\) (the original extracted
convention is defective and is preserved in the annex); its potential is
\(u/2\). These positive rescalings preserve comparison, full mass and
integrability. Smooth backgrounds and the actual currents are unchanged.
No Xia big-class corollary is applied to the boundary class.

Fix a Kahler form \(\Omega\) and write \[
M_k(\psi,\Omega)=\int_X\langle P^k\rangle\Omega^{n-k},\quad M_0=\int_X\Omega^n>0,
\quad d=\max\{k:M_k>0\},\quad m=n-d.
\] Then \(d<n\), and \(d\) is independent of \(\Omega\) and of an anchor
in \(\mathcal{C}_\psi\). Indeed
\(P+\varepsilon\Omega\geq\varepsilon\Omega\) makes
\(\alpha+\varepsilon[\Omega]\) big, and positive multilinearity and the
upper mass bound give \[
0\leq\sum_{k=0}^n\binom nk\varepsilon^{n-k}M_k
=\int_X\langle(P+\varepsilon\Omega)^n\rangle
\leq\operatorname{vol}(\alpha+\varepsilon[\Omega])\longrightarrow0.
\] Thus \(M_n=0\). For any other Kahler form \(\omega\), compactness
gives \(c\Omega\leq\omega\leq C\Omega\). Expanding the positive power
differences gives, for every positive \((k,k)\) product current \(S\),
\[
c^{n-k}\int S\Omega^{n-k}\leq\int S\omega^{n-k}\leq C^{n-k}\int S\Omega^{n-k}.
\] This proves invariance of the positivity tests. DDNL Proposition 2.1
with \(k\) reference slots proves the anchor assertion. For
more-singular \(u_1,\ldots,u_k\), DDNL Theorem 2.4 gives \[
0\leq\int\langle T_{u_1}\cdots T_{u_k}\rangle\Omega^{n-k}\leq M_k.
\] Consequently every such product with \(k>d\) is zero as a current. A
positive current of zero Kahler trace mass is zero: in a local frame,
positivity bounds the absolute coefficient measures by a constant times
the sum of diagonal trace measures; a finite covering makes this global.
Equivalently, normalize its coefficient cone to unit trace; its compact
slice bounds every coordinate, so trace zero forces all coefficients to
vanish. This argument applies also to cluster differences later. No
identification with class numerical dimension is made.

Define initially for a Kahler form \(\omega\) \[
\mathcal{H}_{\psi,\omega}=\left\{u\in\operatorname{PSH}(X,\theta):u\leq\psi+C,
\quad\int\langle T_u^d\rangle\omega^m=\int\langle P^d\rangle\omega^m\right\}.
\] \textbf{Common-domain theorem.} This set is independent of the full
auxiliary form, including its class. Denote it \(\mathcal{H}_\psi\).

For \(d\geq1\), suppose equality holds for \(\omega_0\). Choose \(C>0\)
so that \(D=C\omega_0-\omega_1\) is Kahler. For \(0\leq j\leq m\) define
the scalar deficit \[
\Delta_j=\int_X(\langle P^d\rangle-\langle T_u^d\rangle)\omega_1^jD^{m-j}.
\] It is nonnegative by DDNL Theorem 2.4, applied with \(d\) copies of
\((\theta,\psi)\) versus \((\theta,u)\), \(j\) copies of
\((\omega_1,0)\), and \(m-j\) copies of \((D,0)\). Each background is
smooth closed real and pseudoeffective, and \(\psi\) is less singular
than \(u\). All products are finite by BEGZ. Expansion of the smooth
complementary power yields \[
0=C^m\int_X(\langle P^d\rangle-\langle T_u^d\rangle)\omega_0^m
=\sum_{j=0}^m\binom mj\Delta_j.
\] Every coefficient is strictly positive; each \(\Delta_j\) is
therefore zero. The term \(j=m\) gives equality for \(\omega_1\).
Interchanging the two forms proves equality of domains. This reasoning
never asserts positivity of the difference of product currents. For
\(d=0\) the two empty products are both \(1\), so every deficit is
identically zero and the condition is automatic. The common degree-zero
domain is simply all more-singular \(\theta\)-psh potentials.

For \(d\geq1\) define the joint-linear domain \[
\mathcal{F}^1_\psi=\left\{u\in\mathcal{H}_\psi:
\sum_{a=0}^d\binom da\int_X|u-\psi|\langle T_u^aP^{d-a}\rangle\omega^m<\infty\right\}.
\] For \(d=0\) use \(\int|u-\psi|\omega^n<\infty\). This definition is
independent of the full form as well. In fact for each positive mixed
current \(Q_a\) the pointwise comparison of \(\omega_0,\omega_1\) gives
\[
c^m\int|u-\psi|Q_a\omega_0^m\leq\int|u-\psi|Q_a\omega_1^m
\leq C^m\int|u-\psi|Q_a\omega_0^m.
\] These inequalities hold also for infinite integrals. There are only
finitely many positive summands. For \(d=0\), both quasi-psh functions
are locally integrable against a smooth volume on compact \(X\), so this
finiteness is automatic.

Subtract a constant from \(u\) so that \(u\leq\psi\), and simultaneously
subtract a constant from the reference and potentials if a negative
normalization is needed. The change in each absolute-weight integral is
at most a constant times a finite mixed mass. In Vu's source convention
the diagonal joint linear energy is exactly \[
\frac d2\sum_{a=0}^d\binom da\int_X(\psi-u)\langle T_u^aP^{d-a}\rangle\omega^m.
\] Indeed the sum of \(d\) source potential differences is
\(d(u-\psi)/2\), and there are \(\binom da\) subsets of size \(a\). Thus
the displayed domain agrees exactly with the source diagonal
joint-linear domain for \(d\geq1\), independent of the numerical
weighting of its finitely many positive terms. The source does not
define a zero-slot energy; our \(d=0\) definition is separately
justified. Every member of \(\mathcal{C}_\psi\) belongs to
\(\mathcal{F}^1_\psi\), by bounded differences, BEGZ finiteness and DDNL
mass equality.

If \(\chi\sim\psi\), the dominance conditions are equivalent, \(d\) and
every reference mass agree by DDNL, and hence
\(\mathcal{H}_\chi=\mathcal{H}_\psi\). Vu Corollary 3.2 applies with the
same smooth relative current to identify
\(\mathcal{F}^1_\chi=\mathcal{F}^1_\psi\) when \(d\geq1\). For \(d=0\) it
follows from \(|u-\chi|\leq|u-\psi|+\|\chi-\psi\|_\infty\) off the
common pole set and smooth-volume integrability. This states domain
invariance, not equality of the measures produced by different actual
potentials.

The full-mass domain is convex. For \(d\geq1\) this is Vu Theorem 3.4
with the checked top-degree relative background; for \(d=0\) it is
convexity of psh functions and of the dominance condition. Likewise
\(\mathcal{F}^1_\psi\) is convex by the same theorem, or directly for
\(d=0\). For a finite family \(u_1,\ldots,u_N\in\mathcal{H}_\psi\), its
average \(w\) has full mass. Expanding \(\langle T_w^d\rangle\) shows
its trace mass is the average of \(N^d\) ordered mixed masses. Each is
at most the reference mass by DDNL, all are nonnegative, and their
average equals that mass. Therefore every ordered mixed degree-\(d\)
mass equals the reference mass, against every Kahler \(\omega^m\). This
includes products involving the anchor and repetitions. For \(d=0\) it
is the empty-product identity. This supplies mixed full mass without an
unsupported polarization of scalar diagonal equalities.

The common linear domain can contain genuinely unbounded differences. On
\(X=\mathbb{P}^1_y\times\mathbb{P}^1_z\), put \(\theta=p_1^*\kappa\),
\(\psi=0\) and \(\Omega=p_1^*\kappa+p_2^*\kappa\), with
\(\int\kappa=1\). This is a nef zero-volume class by the cup-product
statement, and \(d=1\), \(M_1=1\). Choose \(0<a<1/2\) and, in a small
disk near \(y=0\), set \(v(y)=-b(-\log|y|)^a\). Continue using a smooth
cutoff on a fixed annulus. Near the puncture the radial second
derivative in \(\log|y|\) is \(ba(1-a)(-\log|y|)^{a-2}>0\). The cutoff
contributes a bounded negative Hessian only on that annulus, so
sufficiently small \(b>0\) ensures \(S=\kappa+dd^cv\geq0\) globally. The
extension takes value minus infinity at the puncture and is subharmonic
there by the distributional radial calculation. Subtract a constant to
have \(v\leq0\).

The radial boundary flux tends to zero, since it is a positive constant
times \(ba(-\log r)^{a-1}\). Thus \(dd^cv\) has no atom at zero;
elsewhere the current is smooth. Stokes gives \(\int S=1\), and deleting
its pole removes no mass. For \(u=p_1^*v\), its degree-one nonpluripolar
mass against \(\Omega\) is \(1\). The integral of \(-v\) against smooth
\(\kappa\) is finite by the exponential decay of radial area after
\(t=-\log r\). Its singular integral against \(dd^cv\) is a positive
constant times \[
\int^\infty t^a t^{a-2}\,dt=\int^\infty t^{2a-2}\,dt<\infty.
\] The remaining annulus is smooth. Pulling back and integrating the
second-factor area therefore proves both degree-one joint integrals
finite. Hence \(u\in\mathcal{F}^1_\psi\), with \(u-\psi\) actually
unbounded. This example is only a witness to strict extension in one
ambient case; the general theorem assumes no fibration.

This proves proposals 1 and 2 in their stated operator-domain scope. No
moving unbounded-weight perturbation limit or energy minimum-reference
theorem follows from these domain statements.

\section{n2. Bare-current reconstruction and mixed decreasing
continuity}\label{n2.-bare-current-reconstruction-and-mixed-decreasing-continuity}

This note uses n1. The new assertion is convergence of the product
current itself, not only of one trace measure.

We use Vu, \emph{Relative non-pluripolar product of currents}, Ann.
Global Anal. Geom. 60 (2021), 269--311, arXiv:2004.11111v1, Lemma
4.1(i): on \(U\subset\mathbb{C}^n\), psh factors \(p_{i,l}\geq p_i\)
converging in \(L^1_{\mathrm{loc}}\), and closed positive relative
backgrounds \(R_l\to R\) satisfying Condition (*), give \[
\liminf_l\int_U\left\langle\bigwedge_{i=1}^d dd^cp_{i,l}\mathbin{\dot\wedge}R_l\right\rangle\Phi
\geq\int_U\left\langle\bigwedge_{i=1}^d dd^cp_i\mathbin{\dot\wedge}R\right\rangle\Phi
\] for every compactly supported smooth weakly positive complementary
test form \(\Phi\). The number of product slots here is \(d\), distinct
from \(m=n-d\) in n1. We use only \(R_l=R=1\), a closed positive
\((0,0)\) current. The relative product is then exactly the ordinary
nonpluripolar product (Vu Definition 3.3). The printed Condition (*) has
an overloaded index in its power, preserved literally in the annex. At
relative degree zero its displayed power is the empty product. Locally
uniform mass and \((q_l-q)1\to0\) follow from local integrability and
dominated convergence on compact subsets for psh \(q_l\downarrow q\)
with bounded limit. The source also explicitly lists constant sequences
as examples. No variable singular relative background is used.

For additional transparency, here is a complete reconstruction of the
required constant-background \emph{decreasing} case. It avoids relying
on a transfer from top measures to currents and keeps every local
product factor bounded before using bounded convergence. The convergence
foundation is DDNL Proposition 2.2: on an open set in \(\mathbb{C}^n\),
uniformly bounded psh factors converging in capacity in every slot,
multiplied by uniformly bounded quasicontinuous weights converging in
capacity, give weak convergence of their weighted Bedford--Taylor top
products. Quasi-psh functions are quasicontinuous and decreasing
sequences with a fixed background converge in capacity, with the
qualifications explained next.

Precisely, Guedj--Zeriahi, \emph{Intrinsic capacities on compact Kahler
manifolds}, J. Geom. Anal. 15 (2005), 607--639, Proposition 3.7 assumes
bounded potentials: decreasing bounded \(A\Omega\)-psh functions
converge in \(\operatorname{Cap}_{A\Omega}\). For arbitrary negative
quasi-psh \(v_l\downarrow v\not\equiv-\infty\) in a fixed smooth
background, choose \(A\Omega\) dominating its lower Hessian bound. The
bounded truncations \(\max(v_l,-K)\) decrease to \(\max(v,-K)\); outside
\(\{v\leq-K\}\) the convergence tests for the truncations and original
functions coincide. Proposition 3.6 gives \[
\operatorname{Cap}_{A\Omega}(\{v<-t\})\leq
\frac{1}{t}\left(\int_X(-v)(A\Omega)^n+n\int_X(A\Omega)^n\right).
\] The set with \(\leq-K\) is contained in \(\{v<-K+1\}\) for \(K>1\),
so its capacity tends to zero. First take \(l\to\infty\) for fixed
\(K\), then \(K\to\infty\). This proves the needed general decreasing
convergence. Proposition 3.10 compares this global capacity with local
Bedford--Taylor capacities on relatively compact coordinate pieces.
Corollary 3.8 gives quasicontinuity. Every function in the present
application is global \(\theta\)-psh plus a smooth local background;
bounded local truncations preserve capacity convergence and
quasicontinuity. These applications use a dominating Kahler form, not
capacity in the pseudoeffective background itself. The source's bounded
hypothesis is never dropped silently.

Let \(p_{i,l}\downarrow p_i\) be these local psh representatives, on a
chart containing the support of \(\Phi\). Subtract constants separately
in each slot so that all \(p_{i,l}\leq0\) on a slightly smaller chart;
local upper bounds follow from the fixed first term. Put \[
s=\sum_{i=1}^d p_i,\qquad b_K=\min(1,\max(s+K,0)),\qquad p^K_{i,l}=\max(p_{i,l},-K).
\] For \(d\geq1\), \(0\leq b_K\leq1\), and its nonzero locus is
\(\{s>-K\}\subset\bigcap_i\{p_i>-K\}\). Since \(p_{i,l}\geq p_i\), all
factors, including the limit factors, equal their bounded truncations on
that nonzero locus. This is a statement about the nonzero locus, not
necessarily the closed support. At the boundary the weight is zero. The
nonzero locus is plurifinely open. BEGZ locality consequently identifies
the weighted nonpluripolar products there with weighted bounded
Bedford--Taylor products. The functions \(p^K_{i,l}\) are uniformly
bounded psh and converge in capacity. The fixed weight \(b_K\) is
bounded quasicontinuous. Thus \[
\int b_K\left\langle\bigwedge_i dd^cp_{i,l}\right\rangle\Phi
\longrightarrow
\int b_K\left\langle\bigwedge_i dd^cp_i\right\rangle\Phi. \tag{C1}
\] Here DDNL's top-product formulation legitimately covers complementary
smooth forms: express \(\Phi\) in a constant \((m,m)\) basis with smooth
compactly supported scalar coefficients. Each constant basis element is
a linear combination of products of \(m\) constant positive \((1,1)\)
forms. To see the span, elementary \((1,1)\) forms span all such
matrices over \(\mathbb{C}\), and their wedges span all \((m,m)\) basis
monomials; real parts give the real span. Every real Hermitian \((1,1)\)
form is a difference of two positive definite forms after adding a large
identity matrix. Expanding these differences gives the required
positive-factor decomposition. The constant positive factors have smooth
quadratic psh potentials, bounded on the chosen chart. Append them as
the remaining slots and incorporate scalar coefficients into the
quasicontinuous weights. DDNL Proposition 2.2 then proves (C1) term by
term. No unbounded product slot is inserted into that proposition or Vu
v3 Theorem 3.3.

For weakly positive \(\Phi\), positivity yields the lower bound by the
left side of (C1). As \(K\to\infty\), \(b_K\) increases to \(1\) off the
union of the product poles, and the limit current gives no mass to those
poles. Monotone convergence therefore proves the displayed
lower-semicontinuity inequality in exactly the decreasing setting needed
here. This also independently verifies the specialization of Vu Lemma
4.1. For \(d=0\) there are no factors, and all products are identically
\(1\); no cutoff argument is needed.

\textbf{Bare-current theorem.} If \(u_l\downarrow u\) with
\(u_l,u\in\mathcal{H}_\psi\), then \[
\langle T_{u_l}^d\rangle\rightharpoonup\langle T_u^d\rangle
\quad\text{as }(d,d)\text{ currents}. \tag{C2}
\] More generally, for \(u_{i,l}\downarrow u_i\) with every member and
limit in \(\mathcal{H}_\psi\), \(1\leq i\leq d\), \[
\langle T_{u_{1,l}}\cdots T_{u_{d,l}}\rangle
\rightharpoonup\langle T_{u_1}\cdots T_{u_d}\rangle. \tag{C3}
\] Every product in (C3) has the same trace mass \(M_d(\psi,\Omega)\) by
n1. Positivity and the coefficient/trace bounds give uniformly bounded
coefficient measures on compact \(X\). Weak compactness of these
measures gives a cluster current \(S\). This compactness can be seen
directly by applying weak-* compactness to each coefficient on finitely
many charts and using a diagonal subsequence; distributional closedness
and positivity pass to the limit. On each chart take a smooth real local
potential \(\rho\) for \(\theta\), which need not be psh. The functions
\(p_{i,l}=\rho+u_{i,l}\) are psh, decrease to \(p_i=\rho+u_i\), and
converge in \(L^1_{\mathrm{loc}}\): off the null pole sets, \[
0\leq p_{i,l}-p_i\leq p_{i,1}-p_i,
\] whose right side is locally integrable. Both limits are not
identically minus infinity on any chart. Lower semicontinuity against
every local weakly positive \((m,m)\) test shows \[
S-\langle T_{u_1}\cdots T_{u_d}\rangle\geq0.
\] This is a positive-current difference obtained by the cluster
argument, not the false assertion that differences of arbitrary actual
products are positive. Passing to the trace mass along the cluster
subsequence gives mass \(M_d\) for \(S\), equal to that of the target.
Its positive difference has zero trace and is zero by n1. All clusters
coincide, proving (C3), and repetition of slots gives (C2). The argument
uses every complementary test bidegree, not merely \(f\Omega^m\).
Smooth-test convergence and uniform trace bounds also give convergence
against continuous coefficient test forms, by uniform smooth
approximation.

For \(u\in\mathcal{H}_\psi\) subtract a constant to have \(u\leq\psi\),
and define \(u_L=\max(u,\psi-L)\). This is \(\theta\)-psh, lies in
\(\mathcal{C}_\psi\), and decreases to \(u\). DDNL comparison gives full
mass (or directly bounded-type equality does). Thus (C2) proves
bare-current convergence for all real levels \(L\to\infty\), not just an
integer subsequence: a hypothetical failing sequence of levels contains
an increasing subsequence, to which the theorem applies. Restoring
constants changes no product. Every decreasing full-mass approximation
with this limit gives the same current. Simultaneous mixed truncations
are covered by (C3). These assertions require the limit to remain in the
full-mass domain; mass loss would invalidate the cluster uniqueness
step.

For each fixed smooth real complementary form \(\eta\), (C2)--(C3) imply
weak convergence of the signed measures obtained by wedging with
\(\eta\). If additionally \(\eta_l\to\eta\) in \(C^0\) coefficient norm,
the coefficient/trace bound shows \[
\|\langle T_{u_l}^d\rangle\wedge(\eta_l-\eta)\|_{\mathrm{TV}}
\leq c_{X,\Omega,d}M_d\|\eta_l-\eta\|_{C^0}\longrightarrow0.
\] Thus the corresponding moving \emph{smooth, uniformly converging}
factors are allowed. This proves no convergence for arbitrary rough
moving factors or unbounded weights.

\section{n3. Minimum full-form data on the entire common
domain}\label{n3.-minimum-full-form-data-on-the-entire-common-domain}

This note uses n1 and n2. Let \[
W=\operatorname{span}_{\mathbb{R}}\{\omega^m:\omega\text{ is a smooth Kahler form on }X\},
\qquad L_u(\eta)=\langle T_u^d\rangle\wedge\eta,
\qquad N=\bigcap_{u\in\mathcal{C}_\psi}\ker L_u.
\] \(W\) is the algebraic span of actual smooth closed real
complementary forms, not a formal tensor space and not an asserted
identification with every closed \((m,m)\) form. For \(\eta\in W\), its
wedge with the order-zero current is a finite signed measure. This
depends on the actual form \(\eta\), not on a presentation as a sum of
powers. Linearity of ordinary smooth multiplication proves presentation
independence. Set \(Q_W=W/N\) and let \(q_W:W\to Q_W\) be the quotient
map.

\textbf{Kernel extension.} The same \(N\) is the intersection of kernels
over \(\mathcal{F}^1_\psi\) and also over \(\mathcal{H}_\psi\). One
inclusion follows from
\(\mathcal{C}_\psi\subset\mathcal{F}^1_\psi\subset\mathcal{H}_\psi\).
Conversely, if \(\eta\in N\), the bounded-type truncations \(u_L\) of
any \(u\in\mathcal{H}_\psi\) satisfy \(L_{u_L}(\eta)=0\). The
bare-current convergence in n2, with fixed smooth \(\eta\), implies
\(L_u(\eta)=0\). This proves the reverse inclusion. Extending the kernel
to the whole domain is substantive content; merely defining a quotient
would not prove this assertion.

Consequently, for any two Kahler forms, \[
L_u(\omega_0^m)=L_u(\omega_1^m)\text{ for every }u\in\mathcal{H}_\psi
\quad\Longleftrightarrow\quad
q_W(\omega_0^m)=q_W(\omega_1^m). \tag{Q1}
\] The equivalence is unchanged if the tests are restricted to
\(\mathcal{C}_\psi\) or \(\mathcal{F}^1_\psi\). Each \(L_u\) induces a
unique linear measure-valued map \(\overline L_u\) on \(Q_W\). The map
\[
Q_W\longrightarrow\prod_{u\in\mathcal{H}_\psi}\mathcal{M}(X),\qquad
q_W(\eta)\longmapsto(L_u(\eta))_u
\] is injective. Here \(\mathcal{M}(X)\) denotes finite signed measures.
This identifies \(Q_W\) with the image of the actual family evaluation
map; it is an operator-data quotient, not a geometric space carrying the
measures.

The precise universal minimality statement is elementary but useful. If
a linear data map \(r:W\to Z\) is sufficient for the family in the sense
that each \(L_u=B_u\circ r\) with linear \(B_u:r(W)\to\mathcal{M}(X)\),
then \(\ker r\subset N\). The unique linear map \(h:r(W)\to Q_W\)
defined by \(h(r(\eta))=q_W(\eta)\) is well-defined, and
\(q_W=h\circ r\). Conversely the quotient itself is sufficient. Thus
every sufficient linear datum refines \(Q_W\); the coarsest equivalence
relation is exactly equality in \(Q_W\). For data maps restricted to
admissible powers, even without linearity, sufficiency implies that
equal data give equal family outputs, hence equal quotient images by
(Q1); the quotient image therefore factors uniquely through the realized
data image. This is minimality only for this specified operator family
and reference span. It proves no universal optimality for every
conceivable zero-volume theory.

There is also a factorwise kernel characterization. For \(d\geq1\), \[
\eta\in N\quad\Longleftrightarrow\quad
\langle T_{u_1}\cdots T_{u_d}\rangle\wedge\eta=0
\text{ for every }u_1,\ldots,u_d\in\mathcal{H}_\psi. \tag{Q2}
\] To prove it first use bounded-type factors. For \(t_i>0\), the convex
combination \(w=(\sum t_i)^{-1}\sum t_iu_i\) is in \(\mathcal{C}_\psi\),
and \[
0=(\sum t_i)^d L_w(\eta)
=\sum_{i_1,\ldots,i_d}t_{i_1}\cdots t_{i_d}
\langle T_{u_{i_1}}\cdots T_{u_{i_d}}\rangle\wedge\eta.
\] Testing with any smooth scalar function gives a real polynomial zero
on an open positive cone; every coefficient is zero. In particular the
coefficient of \(t_1\cdots t_d\) is \(d!\) times the desired mixed
signed measure. Vanishing against all smooth tests proves its measure is
zero. Simultaneous truncation and mixed bare-current convergence from n2
extend the result to \(\mathcal{H}_\psi\). The reverse direction follows
by repeating all factors. For \(d=0\), \(L_u(\eta)=\eta\) as a smooth
top-degree signed measure, so \(N=\{0\}\) and \(Q_W=W\).

\textbf{Quantitative quotient continuity.} Fix \(\Omega\) and a \(C^0\)
coefficient norm from its metric. There is a constant \(c_{X,\Omega,d}\)
such that for every \(u\in\mathcal{H}_\psi\) and smooth complementary
\(\eta\), \[
\|L_u(\eta)\|_{\mathrm{TV}}\leq c_{X,\Omega,d}M_d(\psi,\Omega)\|\eta\|_{C^0}. \tag{Q3}
\] This follows by the local trace control of coefficient measures in
n1, summing a finite partition of unity. Smooth multiplication is
coefficientwise and has no derivatives. In particular \(N\) is closed in
the \(C^0\) norm restricted to \(W\), as well as in its smooth topology.
Define \[
\rho_\Omega(q_W(\eta))=\inf_{\nu\in N}\|\eta+\nu\|_{C^0}.
\] It is a norm on the algebraic quotient (no completeness asserted). If
it is zero, (Q3) implies \(L_u(\eta)=0\) for every bounded-type \(u\),
so \(\eta\in N\). The remaining norm axioms follow from the infimum
definition. Taking the infimum in (Q3) gives \[
\|\overline L_u(\xi)\|_{\mathrm{TV}}\leq c_{X,\Omega,d}M_d(\psi,\Omega)\rho_\Omega(\xi). \tag{Q4}
\] Equivalent metric norms give equivalent quotient norms. Thus this
temporary metric used to state continuity is not additional operator
reference data. Each induced operator is continuous. If
\(u_l\downarrow u\) in \(\mathcal{H}_\psi\) and \(\xi_l\to\xi\) in this
norm, then
\(\overline L_{u_l}(\xi_l)\rightharpoonup\overline L_u(\xi)\): use (Q4)
on \(\xi_l-\xi\), and n2 on any fixed smooth representative of \(\xi\).
Mixed continuity has the same proof using the common mixed trace mass
and (Q2). These are specified admissible continuities on the quotient;
they involve no unbounded test weights.

\textbf{Raw mass and reference invariance.} For every
\(u\in\mathcal{H}_\psi\) and \(\eta\in W\), \[
\int_X L_u(\eta)=\int_X\langle P^d\rangle\eta=:\mathfrak M_\psi(\eta), \tag{Q5}
\] first for each Kahler power by n1 and then by linearity. It descends
to \(Q_W\), since \(L_\psi\) vanishes on \(N\). On an admissible power
its value is strictly positive and equals \[
M_d(\psi,[\omega])=\{\langle P^d\rangle\}\cdot[\omega]^m>0.
\] Closedness of the product and distributional Stokes prove this
cohomological identity; it is a homogeneous degree-\(m\) polynomial of
the class. This does not make the measure cohomological. Scaling
\(\omega\) by \(a>0\) multiplies its quotient image, operators and
masses by \(a^m\). Adding a constant to \(u\) leaves its current and
operator unchanged. No division by \(\operatorname{vol}(\alpha)\) or by
this positive lower-degree mass is performed.

Changing the reference from \(\psi\) to a bounded-equivalent \(\chi\)
leaves \(\mathcal{C}\), \(\mathcal{H}\), \(\mathcal{F}^1\), \(d\), \(W\),
\(N\), and \(Q_W\) unchanged by n1 and the kernel construction. It also
leaves (Q5) unchanged, by equality of full-mass integrals for every
power and linearity. The reference potentials themselves can produce
different measures; the claim is equality of the family and its datum,
not \(L_\psi=L_\chi\). An energy anchor is unnecessary for this
operator; if one is separately chosen, it fixes an energy origin and is
not measure-carrier data.

A gauge change \(\theta' =\theta+dd^cg\), \(u'=u-g\), \(\psi'=\psi-g\),
with \(g\) smooth real, preserves every actual current, comparison and
difference. It preserves all domains, operators and kernel data under
this identification. Any two smooth real representatives of \(\alpha\)
on compact Kahler \(X\) are related in this way by the smooth \(dd^c\)
lemma: Deligne--Griffiths--Morgan--Sullivan, \emph{Real homotopy theory
of Kahler manifolds}, Invent. Math. 29 (1975), 245--274, Lemma 5.11 and
Remark 5.14. Exactly, a smooth \(d\)- and \(d^c\)-closed \(d\)-exact
form is \(dd^c\)-exact, also over \(\mathbb{R}\). The difference of these
representatives is real closed of pure type \((1,1)\), hence
\(d^c\)-closed and exact; its degree-two primitive is a smooth real
function after rescaling the convention. The complete authenticated
printed proof is in the annex. Temporary local potential gauges likewise
have no effect on products by BEGZ.

\textbf{An explicit geometric kernel in a test case.} Let
\(X=\mathbb{P}^1_y\times\mathbb{P}^1_z\), let \(\kappa\) be normalized
Fubini--Study area with integral \(1\), and put \(A=p_1^*\kappa\),
\(B=p_2^*\kappa\), \(\theta=A\), \(\psi=0\). The class is nef,
\(\alpha^2=0\), and therefore has volume zero by the nef cup-product
statement in n1. Here \(d=1\), \(m=1\). In this case \(W\) is every
smooth closed real \((1,1)\) form: given such \(\eta\), both
\(\eta+C(A+B)\) and \(C(A+B)\) are Kahler for sufficiently large \(C\),
and their difference is \(\eta\).

For every actual positive current \(T\in[A]\), the ordinary
smooth-factor measure \(T\wedge A\) has mass \([A]^2=0\) by Stokes and
is positive, so it is zero. In a product coordinate chart write the
positive coefficient matrix of \(T\) in \(y,z\) coordinates. The
\(z\bar z\) diagonal measure is zero, because it is the coefficient
detected by wedging with the positive horizontal form \(A\). Positivity
forces the off-diagonal measures to vanish as well: relative to the sum
of diagonal measures, the coefficient matrix has positive semidefinite
Radon--Nikodym density almost everywhere, and
\(|h_{12}|^2\leq h_{11}h_{22}=0\). Thus \(T\) has only a horizontal
\(y\bar y\) coefficient. Its degree-one nonpluripolar part has the same
property.

It follows that the kernel is exactly \[
N=\{\eta\in W:A\wedge\eta=0\}. \tag{Q6}
\] Necessity follows by testing \(u=\psi=0\). Conversely
\(A\wedge\eta=0\) means the smooth vertical coefficient of \(\eta\) is
zero pointwise. A purely horizontal current then wedges to zero with
\(\eta\), proving sufficiency for every actual potential, in particular
for the full domain. This is an explicit factorwise geometric
characterization in this example, not a general fibration assumption.

The quotient in (Q6) is infinite-dimensional. Choose countably many
disjoint coordinate disks in \(\mathbb{P}^1_z\) with an open region left
over, and nonzero smooth real bumps \(g_j\) supported in those disks.
The exact forms \(\eta_j=p_2^*dd^cg_j\) are in \(W\). A finite linear
combination belongs to \(N\) only if \(dd^c(\sum c_jg_j)=0\) on
\(\mathbb{P}^1\). Such a smooth real function is harmonic and constant:
integration by parts gives \(\int dh\wedge d^ch=0\), so its differential
is zero. The function vanishes on the remaining open region, so the
constant is zero; disjoint supports imply each \(c_j=0\). Thus the
quotient images of these exact forms are linearly independent, despite
every \(\eta_j\) having zero cohomology class and zero scalar mass.

In particular, for a nonconstant smooth \(g(z)\) and sufficiently small
nonzero \(t\), \[
\omega_0=A+B,\qquad\omega_1=A+B+t\,p_2^*dd^cg
\] are Kahler and cohomologous, but \[
L_0(\omega_1)-L_0(\omega_0)=tA\wedge p_2^*dd^cg\ne0,
\qquad\int L_0(\omega_0)=\int L_0(\omega_1)=1.
\] This reconstructs the same-class obstruction analytically and
strengthens it by the exact kernel and infinite-dimensional quotient
calculation.

The actual carrier of every measure remains the \(n\)-dimensional
manifold \(X\); the bare product has bidegree \((d,d)\). Neither \(Q_W\)
nor its admissible-power image is a geometric \(d\)-dimensional carrier.
The minimum datum here is the quotient image of the full complementary
form, for fixed \(X\), actual bounded reference type and actual input
current, together with conventions. No numerical-dimension theorem or
independently geometric general classification has been proved.

\section{n4. Perturbation limits, obstruction and residual
scope}\label{n4.-perturbation-limits-obstruction-and-residual-scope}

This note uses n1--n3. There are two distinct verified perturbation
statements. Neither permits arbitrary unbounded weights.

\textbf{Decreasing full-mass paths in the original background.} Let
\(\varepsilon_j\downarrow0\), let \(u_j\downarrow u\) with every
\(u_j,u\in\mathcal{H}_\psi\) and \(u_j\in\operatorname{PSH}(X,\theta)\),
and let \(\omega_j\) be smooth Kahler forms with
\([\omega_j]\to\lambda\) for a Kahler class \(\lambda\). Suppose there
is a smooth Kahler form \(\omega\) such that \[
\rho_\Omega\big(q_W(\omega_j^m)-q_W(\omega^m)\big)\longrightarrow0. \tag{P1}
\] The form \(\omega\) need not represent \(\lambda\): both conditions
are stated separately, since one controls scalar lower-degree
coefficients and the other controls the operator reference. Then \[
\frac{\langle(T_{u_j}+\varepsilon_j\omega_j)^n\rangle}
{\binom nd\varepsilon_j^m}
\rightharpoonup\overline L_u(q_W(\omega^m))=\langle T_u^d\rangle\omega^m. \tag{P2}
\] To prove this expand actual positive currents, and use the
higher-degree vanishing in n1: \[
\frac{\langle(T_{u_j}+\varepsilon_j\omega_j)^n\rangle}
{\binom nd\varepsilon_j^m}
=L_{u_j}(\omega_j^m)+
\sum_{k<d}\frac{\binom nk}{\binom nd}\varepsilon_j^{d-k}
\langle T_{u_j}^k\rangle\omega_j^{n-k}. \tag{P3}
\] Every summand of the remainder is positive. DDNL comparison gives its
unscaled mass at most \(M_k(\psi,[\omega_j])\). Because
\(\langle P^k\rangle\) is closed, this reference mass equals its
cohomology pairing with \([\omega_j]^{n-k}\), a fixed continuous
polynomial. Convergence of classes bounds all these finitely many
masses. Thus the remainder has total variation \(O(\varepsilon_j)\) for
\(d\geq1\), uniformly without pointwise bounds on representatives; for
\(d=0\) there is no remainder. Joint quotient continuity in n3, using
(P1) and the decreasing bare-current theorem, gives the leading limit.
This proves (P2).

Fixed \(u\) in the full domain is allowed by the constant decreasing
sequence. In particular fixed-form perturbation convergence holds for
every such unbounded \(u\), with a positive raw mass. Smooth convergence
\(\omega_j\to\omega>0\) implies both class convergence and (P1), and is
a useful sufficient special case. Equality of quotient images at every
\(j\), together with convergent classes, is another sufficient case even
without convergence of representatives. The theorem supplies
independence for all paths satisfying its exact hypotheses and the same
limiting datum. It does not assume \(u_j\) has bounded differences,
finite linear energy or an \(o(\varepsilon_j)\) rate. Its essential
added hypothesis is that these decreasing currents are positive in the
original background and all inputs and their limit have full
degree-\(d\) mass. It makes no assertion for arbitrary potentials psh
only in \(\theta+\varepsilon_j\omega_j\).

\textbf{Bounded-type perturbed backgrounds.} For completeness,
reconstruct the distinct sufficient rate for that latter setting. Fix
\(v\in\mathcal{C}_\psi\), let \(\omega_\varepsilon\to\omega>0\) smoothly,
and suppose
\(v_\varepsilon\in\operatorname{PSH}(X,\theta+\varepsilon\omega_\varepsilon)\)
has the same actual bounded type as \(\psi\) and \[
\delta_\varepsilon=\|v_\varepsilon-v\|_\infty=o(\varepsilon)
\] off the common pole set. Set \(A=T_v+\varepsilon\omega_\varepsilon\),
\(B=\theta+\varepsilon\omega_\varepsilon+dd^cv_\varepsilon\),
\(h=v_\varepsilon-v\). Positive multilinearity gives the signed top
difference as the sum of \(dd^ch\) insertions against \[
R_{a,\varepsilon}=\langle B^aA^{n-1-a}\rangle,\qquad 0\leq a\leq n-1.
\] These are closed positive \((n-1,n-1)\) currents with no pluripolar
mass, or the smooth current \(1\) when \(n=1\).

The exact literature IBP used here is Vu v3 Lemmas 2.5--2.6 and Theorem
2.7: for a closed positive \((n-1,n-1)\) current \(R\) and bounded
\(R\)-admissible dsh functions \(h,f\), \[
\int f\langle dd^ch\mathbin{\dot\wedge}R\rangle
=\int h\langle dd^cf\mathbin{\dot\wedge}R\rangle
=-\int\langle dh\wedge d^cf\mathbin{\dot\wedge}R\rangle.
\] Admissible means a representation as differences of quasi-psh
functions whose polar sets have zero \(R\)-mass. The bounded difference
\(h\) and any smooth real test \(f\) meet this condition, because the
relative currents have no pluripolar mass. Vu Proposition 2.1(iii)
identifies the signed insertion with the difference of the ordinary
nonpluripolar products; smooth \(dd^cf\) is ordinary multiplication.
Source potentials are divided by \(2\); source \(d^c_s=2d^c\). Both
source IBP sides are half ours, so multiplication by \(2\) gives the
formula displayed here. The correct intrinsic gradient is \[
dh\wedge d^ch\wedge R=\tfrac12dd^c(h^2)\wedge R-h\,dd^ch\wedge R.
\] The printed missing half factors on Vu pages 6 and 11 and the older
Cauchy--Schwarz defect remain literal in the annex and are not
definitions used here. For clarity the cutoff boundary estimate can be
checked with the correct formula: if \(h=\varphi_1-\varphi_2\),
\(f=\varphi_3-\varphi_4\) with negative quasi-psh representatives,
\(dd^c\varphi_i\geq-\Gamma\), put
\(b_k=1+k^{-1}\max(\sum\varphi_i,-k)\). It is bounded between \(0\) and
\(1\), \(dd^cb_k\geq-4\Gamma/k\), and bounded Bedford--Taylor calculus
and Stokes give \[
0\leq\int db_k\wedge d^cb_k\wedge R
=-\int b_kdd^cb_k\wedge R\leq\frac4k\int\Gamma\wedge R.
\] The truncated differences have uniformly bounded sup norms and
diagonal gradient mass at most \(2\|h_k\|_\infty\int\Gamma R\), and
similarly for \(f_k\). Cauchy--Schwarz makes every boundary term with
\(db_k\) vanish. The nonzero cutoff locus, not its closed support, lies
in the simultaneous strict truncation locus. Locality and admissible
nonpluripolar gradient finiteness (the cited Lemmas 2.5--2.6) let the
intrinsic products pass to the limit off the representation poles;
dominated convergence applies since \(R\) and the inserted positive
products give them no mass. Thus the corrected intrinsic formula gives
exactly the scoped IBP application. No erroneous doubled square
expansion is used.

Choose \(C_f\) uniformly with
\(-C_f\omega_\varepsilon\leq dd^cf\leq C_f\omega_\varepsilon\), possible
by smooth convergence to a positive form. IBP gives \[
\left|\int f(\langle B^n\rangle-\langle A^n\rangle)\right|
\leq nC_f\delta_\varepsilon
\sum_{k=0}^d\binom{n-1}k\varepsilon^{n-1-k}M_k(\psi,[\omega_\varepsilon]). \tag{P4}
\] Indeed DDNL Proposition 2.1 applies to the \(n-1\) bounded-equivalent
slots in the big class \(\alpha+\varepsilon[\omega_\varepsilon]\) and
one \((\omega_\varepsilon,0)\) slot. Replacing all the varying potential
slots by \(\psi\) makes their common mixed mass
\(\int\langle(P+\varepsilon\omega_\varepsilon)^{n-1}\rangle\omega_\varepsilon\),
whose positive expansion is precisely the sum in (P4). Since \(d<n\),
all indices are legal; when \(n=1,d=0\) it is the empty-product mass
\(M_0\). The sum is uniformly \(O(\varepsilon^{m-1})\). After rescaling,
the test error is \(O_f(\delta_\varepsilon/\varepsilon)\to0\). The
rescaled masses are uniformly bounded by the same bounded-type
comparison and actual reference polynomial. Smooth tests are dense in
continuous functions, so convergence on smooth tests proves weak measure
convergence. Fixed-current expansion and smooth multiplication then give
the limit \(L_v(\omega^m)\). This is a sufficient rate only; no
optimality or rate-free assertion for perturbed backgrounds is made.

\textbf{An analytic obstruction to dropping full-form data.} In the
product example of n3 use the fixed potential \(u=0\), and alternate the
cohomologous forms \(\omega_0,\omega_1\) along \(\varepsilon_j=1/j\),
choosing \(\omega_0\) for even \(j\) and \(\omega_1\) for odd \(j\). All
backgrounds are Kahler perturbations, all auxiliary classes equal one
class, all potentials are identical and bounded, and the class \([A]\)
is pseudoeffective of volume zero. Direct smooth expansion gives \[
\frac{(A+\varepsilon_j\omega_j)^2}{2\varepsilon_j}
=A\wedge\omega_j+\frac{\varepsilon_j}{2}\omega_j^2.
\] The last term tends to zero in total variation since this is a finite
set of smooth forms. The even and odd subsequences tend to \(A\wedge B\)
and \(A\wedge B+tA\wedge p_2^*dd^cg\), respectively, which are distinct
positive mass-one measures by n3. Thus class convergence alone, even
with zero potential error, does not guarantee a unique limiting operator
on \(X\). The failed hypothesis in (P1) is convergence of the actual
reference datum. This counterexample refutes only a class-only
prescription for this full-form operator; it is not an impossibility
proof for every possible theory sought in the original atom.

\textbf{Catalogue comparison.} New content relative to the supplied
narrower operator result is: the full-mass domain is common to every
auxiliary form and class; its joint-linear subdomain is likewise common;
truncation and mixed decreasing convergence hold for bare currents; the
bounded-family reference kernel extends unchanged to the whole full-mass
domain; a quantitative quotient norm supplies admissible joint
continuity; and an explicit example computes the kernel and its
infinite-dimensional quotient. The quotient-path limit (P2) also covers
decreasing unbounded full-mass inputs positive in the original
background. These are rigorous scoped applications with proofs above,
not literature-priority claims or independent certifications. The
universal quotient definition itself is a bookkeeping characterization,
and is not claimed as an independently geometric classification.

The retained operator data are \(X\) and its complex structure, the
cohomology class with an actual bounded singularity type dominating the
domain, the actual positive input current, the actual-current degree
\(d\), and \(q_W(\omega^m)\) for an admissible full-form reference. The
smooth \(\theta\) is a gauge; a particular bounded-equivalent anchor is
not additional operator datum. The carrier is explicitly \(X\), of
complex dimension \(n\); the product current has bidegree \((d,d)\), raw
positive mass is (Q5), and there is no probability normalization. The
degree need not equal numerical dimension: in the product class \([A]\),
the logarithmic actual reference with current
\([\{y=0\}\times\mathbb{P}^1]\) has zero degree-one nonpluripolar part
and hence \(d=0\), while the smooth reference has \(d=1\). The
logarithmic potential is \(p_1^*\log(|y|^2/(1+|y|^2))\); in our
convention \(dd^c\log|y|^2=\delta_0\) by boundary flux \(4\pi\) and
\(dd^c=(\Delta/4\pi)\) times planar area. Subtracting
\(dd^c\log(1+|y|^2)=\kappa\) gives the asserted divisor current. Its
nonpluripolar part is zero because it is supported on its pole set. The
nef class has numerical dimension \(1\) by its nonzero first cup product
and zero square.

Original P15.2 remains open in the dossier's residual scope. The measure
construction uses a prescribed-type relative full-mass domain and
specified perturbation classes; it supplies no intrinsic geometric
\(d\)-dimensional carrier, no independently geometric general
characterization of the reference kernel, no numerical-dimension
identification, no identification of actual reference mass with a
leading class-volume coefficient, and no independence for arbitrary
genuinely perturbed-background unbounded paths. Full-mass limits and the
reference-datum convergence hypotheses in (P2), and bounded type plus
smooth auxiliary convergence and \(o(\varepsilon)\) error in (P4), must
be retained. No increasing-sequence continuity or varying
unbounded-weight calculus is proved. Proposal 3 about an energy quotient
is deliberately not investigated in this P15.2-only assignment; no P15.1
assertion is needed. The optional degree-zero rigidity question is left
unanswered without treating it as a blocker. Status is partial and
candidate is false.

\section*{Review and supporting evidence}

The complete proof passed an independent Xean native Codex review using
\texttt{gpt-6.1-sol} with \texttt{xhigh} reasoning. Its entire report,
11 premises and 26 source passages were read and bound
to the frozen proof, TeX source and author status. All new mathematical
arguments, actual canonical notes and the author manuscript were read in
full. All 24 selected literal primary-source bodies were authenticated;
unchanged complete primary reads were explicitly reused. The separate
StonyBrook source authenticates the DGMS real smooth $dd^c$ lemma and its
complete selected proof. The GZ bounded capacity hypothesis and Vu bounded
product-factor hypothesis are retained through the proved truncation
applications. Literal source defects are preserved in the dossier; the
new proofs use the correct intrinsic one-half gradient identity.

A fresh copied build reproduced byte-identical mathematical proof,
TeX source and literal primary annex. The full argument differs only in
the execution paragraph containing the newly generated PDF digest; that
single substitution was verified. Execution checks are reproducibility
evidence, not analytic correctness.

Publication edits update review provenance and typesetting and remove
execution bookkeeping. They do not repair a mathematical gap or change
the theorem. The full frozen proof, complete primary units, source audits,
canonical notes, review and root read records accompany this manuscript.
Native editor compilation and local PDF export are recorded separately.
Same-model independent review, compilation and file integrity are not
formal verification or literature-priority certification. Original P15.2
closure requires a separate audit against the exact original question;
the present acceptance is limited to the theorem stated in this paper.

\end{document}
